\chapter{Căn Bệnh Con Nhà Người Ta}
\markboth{Căn Bệnh Con Nhà Người Ta}{Căn Bệnh Con Nhà Người Ta}

\begin{truyen}{Bạn Bằng Tuổi Người Ta}{Comparison Pressure}
\chuhoa{M}{ột bạn ngồi cuối lớp than nhỏ: ``Ở nhà, hễ em được 8 là ba mẹ hỏi sao không được 9 như con cô kia. Em mệt nhất không phải vì học, mà vì lúc nào cũng bị đem so với một đứa hoàn hảo tưởng tượng nào đó.''}

\teacherpair{Thầy nói chậm: ``Nhiều cha mẹ tưởng so sánh là động lực. Thật ra dùng hoài thì nó thành dao mòn: không chém đứt ngay, nhưng cứa lòng tự trọng của con mỗi ngày. Người bị so liên tục sẽ không còn học vì tiến bộ, mà học để đỡ bị chê.''}{The teacher responds: ``Many parents think comparison creates motivation. Used repeatedly, it becomes a dull knife: it may not cut deep at once, but it scrapes away a child's self-worth every day. A constantly compared child no longer studies to improve, but to avoid humiliation.''}

Bạn ấy hỏi: ``Vậy em phải làm sao khi nói thì ba mẹ bảo đó là thương cho roi cho vọt?''

\teacherpair{Em không cần cãi bằng giận dữ. Em cần nói bằng dữ kiện: con đang cố ở đâu, con mệt chỗ nào, và con cần hỗ trợ kiểu gì. Còn nếu người lớn chưa nghe, em càng phải giữ một điều: đừng biến lời so sánh của họ thành giọng nói tự chê bai ở trong đầu em.}{``You do not need to fight with anger. Speak with facts: where you are trying, where you are struggling, and what kind of support you need. And if adults still do not hear you, keep one thing safe: do not let their comparisons become your own inner voice against yourself.''}

Thầy nói thêm, như để trả lại một ít công bằng cho những người làm cha mẹ: ``Nhiều phụ huynh cũng bị đời so với người khác suốt. Họ mang cái đau đó về nhà mà không biết. Thành ra họ không hẳn muốn ác với con; họ chỉ đang lặp lại đúng thứ đã làm mình tổn thương. Nhưng lặp lại nỗi đau của mình lên con thì vẫn là sai.''

\textit{(A student speaks about comparison at home. The teacher names it as slow damage to self-worth rather than healthy motivation.)}
\end{truyen}

\begin{baihoc}
So sánh có thể đẩy một đứa trẻ chạy nhanh hơn trong ngắn hạn, nhưng thường làm nó mệt và ghét chính mình về dài hạn. Một gia đình trưởng thành không nuôi con bằng bảng xếp hạng, mà bằng sự tiến bộ có thật của chính đứa trẻ đó.
\textit{(Comparison may make a child run faster in the short term, but often leaves them exhausted and self-rejecting in the long term. A mature family does not raise a child by rankings, but by the child's real growth.)}
\end{baihoc}

\begin{ghinhoanh}
\textbf{Short English to remember:}

Comparison can wound motivation.

Growth needs guidance, not humiliation.
\end{ghinhoanh}

\begin{tuhocnhanh}
1. Trong đầu em có đang tồn tại một câu so sánh lặp đi lặp lại không? \textit{(Is there a comparison sentence still repeating inside your mind?)}

2. Nếu được nói một câu thật với cha mẹ, em muốn họ hiểu điều gì nhất? \textit{(If you could say one honest sentence to your parents, what would you most want them to understand?)}
\end{tuhocnhanh}

\ngancach